% 并查集的数据结构：元素1..6，集合{1,2,3,4}（根1）与{5,6}（根5）。
% 父指针：parent = (1,1,1,3,5,5)，根指向自己；大小只在根上有意义：size(1)=4，size(5)=2。
% 左：父指针树，箭头由元素指向父指针；右：两个数组。find(4) 沿 4->3->1 到达根1。
\begin{tikzpicture}[font=\figurefont\fontsize{8}{10}\selectfont,>=Stealth,
  el/.style={circle,draw=FigureStroke,line width=.8pt,minimum size=5.8mm,inner sep=0pt,text=FigureInk,fill=white},
  rt/.style={el,fill=FigureBlueFill},
  par/.style={draw=FigureStroke,line width=.9pt,->},
  pathp/.style={draw=FigureYellow,line width=1.6pt,->},
  cell/.style={draw=FigureStroke,line width=.6pt,minimum width=6.4mm,minimum height=5.6mm,inner sep=0pt,text=FigureInk},
  rcell/.style={cell,fill=FigureBlueFill},
  hd/.style={font=\FigureNoteFont\fontsize{7.5}{9}\selectfont,text=FigureMuted},
  note/.style={font=\FigureNoteFont\fontsize{7.5}{9}\selectfont,text=FigureMuted,align=center}]
% 左：两棵树
\node[rt] (n1) at (0.9,0) {1};
\node[el] (n2) at (0,-1.05) {2}; \node[el] (n3) at (0.9,-1.05) {3};
\node[el] (n4) at (0.9,-2.1) {4};
\draw[par] (n2)--(n1); \draw[par] (n4)--(n3);
\draw[pathp] (n3)--(n1);
\draw[pathp] (n4) to[bend right=0] (n3);
\draw[par,->] (n1) edge[loop above,looseness=6,min distance=5mm] (n1);
\node[rt] (n5) at (2.3,0) {5}; \node[el] (n6) at (2.3,-1.05) {6};
\draw[par] (n6)--(n5);
\draw[par,->] (n5) edge[loop above,looseness=6,min distance=5mm] (n5);
\node[note] at (1.15,-2.85) {(a) 父指针树；黄色为$\operatorname{find}(4)$\\经过的路径$4\to3\to1$};
% 右：数组
\begin{scope}[xshift=4.0cm,yshift=-0.35cm]
  \node[hd,anchor=east] at (-0.1,0) {元素$v$};
  \node[hd,anchor=east] at (-0.1,-0.75) {$\operatorname{parent}(v)$};
  \node[hd,anchor=east] at (-0.1,-1.5) {$\operatorname{size}(v)$};
  \foreach \v/\p/\s/\r in {1/1/4/1,2/1/{--}/0,3/1/{--}/0,4/3/{--}/0,5/5/2/1,6/5/{--}/0}{
    \pgfmathsetmacro{\x}{0.32+0.64*(\v-1)}
    \node[hd] at (\x,0) {\v};
    \ifnum\r=1
      \node[rcell] at (\x,-0.75) {\p};
      \node[rcell] at (\x,-1.5) {\s};
    \else
      \node[cell] at (\x,-0.75) {\p};
      \node[cell,text=FigureMuted] at (\x,-1.5) {\s};
    \fi
  }
  \node[note] at (1.92,-2.5) {(b) 实际存储的两个数组；蓝色列是根，\\$\operatorname{parent}$指向自己，$\operatorname{size}$只在根上使用};
\end{scope}
\end{tikzpicture}
